\chapter{Intro}

This projects goal is to create a game engine, similar in design to what is currently available from big commercial engines.
As this is a large undertaking that far exceeds the limits and practicality of a thesis project, I will be focusing on giving the project a good foundation. This means considerable time will be spent investigating solutions to issues other projects do not care about. By the end I hope to have a solid foundation and core systems for the engine that I plan to use and extend in the coming years.
Planned parts:
- CI/CD pipelines
- Choosing a build system that is usable with large projects
- Dynamic DLL loading to allow for hot reloading of the game code
- Basic Entity system
- Reflection system for runtime code metadata

The project is a collaboration between me and a friend of mine, she will be writing the renderer and donating some other bits of code.
The project is hosted on github and is MIT licensed.


\section{What is a game engine?}
A game engine is a highly complex software that allows the creation of computer games. They can vary is what functionality they provide and how the game maker interacts with them. They can have graphical interfaces and easy to understand UI, or be code only and need programming knowledge. The most wildly known games engines (like Unity and UE4) provide a single editor application that the user (the game maker) interacts with. These programs have many built in tools that help with creating a game.

There are game engines witch are availabe for free and are used by many companies for buiilding games. While other games use either a purpuse built engine or one that the comany develoed inhouse.


\section{How games work}
Games are real time graphical applications, meaning that they redraw the contents of the screen fast enough where they form a continous animation. Films work the same way but as they are not interactive they can work with 15 frames per second (FPS), but as games need to react to the user's input they should strive to draw as many frames as the hardware can process. For games developed for PCs the recommended minimum is 60 frames per second. This means the game has to finish redrawing and updating everything in the game world 60 times a second. ~~This gives us **16.7 milliseconds** to do everything.~~

Engines are the core of a games code, they provide many of the utility functions that are needed to draw the game, gather input, load files and many more. Just as there are many different styles of games, there are many different game engines. Some are purpose built engines for a specific game with no intention to be reusable. While others are general purpose engines that can be used to make a wide range of games. These game engines will also warry in what they provide, for example an engine for a single player game will probbaly not have any networking code.

In many cases companies will have their own engine developed for their own games, this is usefull for them as they can modify the core engine code for each game as needed. This can result in higher integration and optimization between the engine and the game. Not many of these so called "in house" engines are available to the public to see. The biggest two that fit this description are Valve's Source engine and Epic Gemes' Unreal Engine, both of these engiens have a long history going back to the 90's and the start of the computer gameing industry.

Because there are countless engines most of witch is not available online I will be only using a few of the most known engines for comparison.
The 3 most popular freely available game engines:
- Unity
- UE4/5
- Godot
On top of these I will be using Valve's Source engine as it is source available and in some cases has simpler implementations of certain tools/functionality

\subsection{Unity}

Unity is a fully multi purpose engine capable of 2D and 3D graphics with physics, and modern rendering features. It's biggest selling point is the large amount of platforms it supports, games made with it can run on any modern PC and every console.

I personally have the most experience with it out of the 3, and it is one of the reasons why I started this project. While using unity, I often found my self wanting to do more complex things than what the engine allowed. This meant that I ended up fighting with the engine more than just using it. This became extremely apparent during my last laboratory project (Ön lab, Unity event system).

\subsection{Unreal 4, 5}

Unreal is one of the oldest engines still under development. Its biggest feature is the high level of visual quality it can provide. Unreal is mostly focused on the AAA style market and FPS / TPS games. It's other big selling point is that it is source available, this means that if you want you can edit the source code and compile a version for your own use (while adheering to the license).

Uneral's problem for my is that it is highly focused on multiplayer FPS games, and some of the settings and facilities it proviedes make no sense outside of that usecase. On top of that it is a large install size witch makes it difficult to use on slower devices.

\subsection{Godot}

Godot is relatively the newest game engine out there, it is developed fully in the open and is MIT licensed. It started out as a 2D game engine but since ( \#re-word version number) it has full support of 3d graphics and is constatnly improving the visual fidelity it can provide.

It has C\# Cpp and their own GScript options for writing your game code this makes complex system and simpler game logic development both easy.

\subsection{Source}

Source is Valve's game engine and was used in almoast all of their games, like Half-Life 2, portal1/2, and many more.

It follows and older style of game engine design where the runtime of the engine is almoast comletly searate from the editing tools. For examle it has a level editor that uses a comletly searate way to draw the designs, and needs to exort it in a format that the runtime can understand.

It has the fastest to use level editor and works


\section{The project}

Base idea
C++ university course, start of the engine
many problems

The idea of a rewrite with others helping

my history:
\begin{itemize}
    \item Unity since 3.4
    \item C\# Cryogenic
    \item C course Relic engine
    \item C++ course ShadowEngineOld
\end{itemize}

The renderer
My friend offered to write the renderer for the project as that is the part that i got most stuck on in the last iteration
For multiple reasons I have decided to make the renderer swappable as to allow the use of a different rendering backend.

The requirements
\begin{itemize}
    \item Modules
    \item Inter module communication
    \item Entities
    \item Editor runtime setup
\end{itemize}

\subsection{Licensing}

The project is currenlty licensed under MIT and will stay that way for the forseeable future
